\section{Reciprocal designs without preferred directions}
\label{sec:isotropic-rigidity}
The exact rigidity is not tied to the four compass directions. We
first identify the property of the displacement law used in the
stopping argument, and then give a rotationally invariant collision
experiment with an intrinsic perimeter normalization. The finite
compass minimax experiment remains the one specified in
Corollary~\ref{cor:stationary-minimax}.

\subsection{A general centered displacement law}
Let $\mu$ be a known probability measure on displacements in $\R^2$
satisfying
\begin{equation}\label{eq:general-command-law}
 \int a\,\mu(\dd a)=0,\qquad
 \sigma_\mu^2=\int|a|^2\,\mu(\dd a)>0,\qquad
 \operatorname{supp}\mu\subset\ell\overline B.
\end{equation}
No symmetry, density, or nondegeneracy of the covariance matrix is
assumed. The obstacles and stationary launch density are as in
Section~\ref{sec:global-response}, with
\begin{equation}\label{eq:general-command-separation}
                              \ell+\Delta<d.
\end{equation}
Draw $Z\sim j$ independently of $a\sim\mu$. The forward attempt
is $(x+Z,a)$, and its reciprocal reverse attempt is
$(x+Z+a,-a)$ with fresh draws from that same nominal joint law.
Let $F_\mu,R_\mu$ denote the pooled means, and put
\[
 g_\mu=F_\mu-R_\mu,\qquad
 T_\mu f(x)=\int f(x+a)\,\mu(\dd a).
\]
The displacement realization is not returned by the sensor.

\begin{theorem}[Design-independent translation rigidity]
\label{thm:general-command-rigidity}
Under \eqref{eq:general-command-law}--\eqref{eq:general-command-separation},
the occupation field satisfies $g_\mu=(T_\mu-I)v$ and is uniquely
determined by $g_\mu$. With
\begin{equation}\label{eq:general-exit-budget}
                     H_\mu=\frac{(D+\Delta+\ell)^2}{\sigma_\mu^2},
\end{equation}
one has
\begin{equation}\label{eq:general-stopping-inverse}
 v(x)=\sup_{\E_x\tau\le H_\mu}
           \E_x\left[-\sum_{k<\tau}g_\mu(X_k)\right],
\end{equation}
where $X$ is the virtual walk with increment law $\mu$ and immediate
stopping is allowed. Consequently
\begin{equation}\label{eq:general-response-periods}
 \operatorname{Per}(F_\mu,R_\mu)
 =\operatorname{Per}(g_\mu)
 =\operatorname{Per}(v)
 =\operatorname{Per}(\mathcal O).
\end{equation}
The same prior-uniform Lipschitz inverse bound holds with constant
$H_\mu$. None of these assertions assumes periodicity of the table.
\end{theorem}
\begin{proof}
The endpoint balance identity applies to each displacement of length
at most $\ell<d$. Independence of $Z$ and $a$ therefore gives
$g_\mu=(T_\mu-I)v$. The positive components are still $C+(-A)$,
with diameter at most $D+\Delta$ and gaps greater than $\ell$.
An exit step from one positive component consequently lands at zero
occupation.

The mean-zero assumption and independence of increments make
$|X_n-x|^2-n\sigma_\mu^2$ a martingale. Before component exit,
and on its exit step, the displacement from $x$ is at most
$D+\Delta+\ell$. Optional stopping first at a bounded time and
then monotone convergence give expected exit time at most $H_\mu$.
This proof uses the trace of the covariance, not its invertibility;
it also applies to a law supported on a line. Markov's inequality and
the Markov property give the same geometric block tail as in
Lemma~\ref{lem:global-exit}.

Stopping the bounded occupation martingale gives a payoff at most
$v(x)$ for every integrable rule. The containing-component exit,
or immediate stopping when $v(x)=0$, attains equality and proves
\eqref{eq:general-stopping-inverse}. Comparing its payoffs for two
forcings gives the Lipschitz bound. The stopping functional commutes
with translations, so every period of $g_\mu$ is a period of $v$.
Such a period permutes the separated expanded components; cancelling
$p_{-A}$ from their laboratory support functions gives a period of
$\mathcal O$. The remaining implications in
\eqref{eq:general-response-periods} follow from translation covariance
and subtraction, exactly as in
Theorem~\ref{thm:global-response-rigidity}.
\end{proof}

For a law with finite lattice support, a sufficiently large finite
killing set gives a stopping polytope as in
Lemma~\ref{lem:stencil-polytope}. For a law with continuous support,
\eqref{eq:general-stopping-inverse} is an exact response-field theorem,
not a claim that finitely many mean samples evaluate the integral
operator exactly. In either case the displacement law is a specified
control, whereas the stationary launch density is unknown.

\subsection{Isotropic collision means and a perimeter deficit}
Now fix a length $t>0$ and take the direction uniformly on the unit
circle:
\begin{equation}\label{eq:isotropic-design}
       a=t(\cos\theta,\sin\theta),\qquad
                      \theta\sim\operatorname{Unif}[0,2\pi).
\end{equation}
Its mean is zero and $\sigma_\mu^2=t^2$. The observation operator
commutes with every orthogonal transformation; there is no distinguished
apparatus axis. Nominal centers and the reciprocal joint law are
still prescribed controls.

Use the translated homothetic supports of
\eqref{eq:registration-model}, with at least two distinct scales, and
assume $2t+\max_i\diam A_i<d_0$. Let $\mathcal L(H)$ denote the
Euclidean perimeter of a planar convex body. From the exact reciprocal
difference at setting $i$, recover the complete expanded-component
collection $\mathcal P_i$ by
Theorem~\ref{thm:general-command-rigidity}. For $P=C+(-A_i)$,
choose a bounded region $E_P$ containing $P+t\overline B$ and no
other component's forward-response support, and put
\begin{equation}\label{eq:isotropic-deficit}
 d_i^{\mathrm{iso}}(P)=\mathcal L(P)
                    -\frac{\pi}{t}\int_{E_P}F_{i,\mu}(x)\dd x.
\end{equation}

\begin{lemma}[Isotropic collision perimeter]
\label{lem:isotropic-perimeter}
For every component and setting,
\begin{equation}\label{eq:isotropic-perimeter}
 \int_{E_P}F_{i,\mu}(x)\dd x=\frac{t}{\pi}\mathcal L(C),
 \qquad d_i^{\mathrm{iso}}(P)=\mathcal L(A_i)>0.
\end{equation}
In particular the deficit is independent of the selected component,
the launch density, and the admissible region $E_P$.
\end{lemma}
\begin{proof}
At a fixed direction $e$, horizontal slicing in that direction gives
\[
 |(C+[-te,0])\setminus C|=t\,w_C(e^\perp),
 \qquad w_C(u)=p_C(u)+p_C(-u).
\]
Tonelli's theorem eliminates the independent launch translation from
the integrated mean. Averaging the strip identity in $e$ gives
\[
 \int_{E_P}F_{i,\mu}
        =\frac{t}{2\pi}\int_0^{2\pi}w_C(n(\theta))\dd\theta
        =\frac{t}{\pi}\mathcal L(C).
\]
The last equality is Cauchy's perimeter formula. To fix its
normalization, for a smooth strictly convex body it follows by
integrating $p_C+p_C''$ over the circle and using periodicity of
$p_C'$. General convex bodies follow by smooth approximation of
support functions and continuity of perimeter. The same formula
shows $\mathcal L(C+H)=\mathcal L(C)+\mathcal L(H)$; these are
classical support-function identities, as in \cite{Schneider}.
Apply them to $P=C+(-A_i)$ and note that reflection preserves
perimeter. The separation assumption justifies using one isolated
response envelope throughout.
\end{proof}

\begin{theorem}[Origin-free rigidity from isotropic pooled responses]
\label{thm:isotropic-registration}
In the isotropic experiment, the exact pooled reciprocal mean pairs
determine the scale ratios, the centered physical first footprint,
and the configurations $\mathcal O-b_i$. No component matching
between settings, preferred direction, homothety origin, periodicity,
finite species list, or boundary smoothness is supplied.

Explicitly, set
\begin{equation}\label{eq:isotropic-ratio}
                   \rho_i=d_i^{\mathrm{iso}}/d_1^{\mathrm{iso}},
\end{equation}
choose $k$ with $\rho_k\ne1$, and form the lower centered support
envelopes $L_i$ in \eqref{eq:registration-envelope}. Then
\begin{equation}\label{eq:isotropic-footprint}
 q=\frac{L_k-L_1}{\rho_k-1}=p_{-A_0},\qquad
                   p_P-\rho_iq=p_{C-b_i}\quad(P\in\mathcal P_i).
\end{equation}
The full geometric ambiguity is exactly
\begin{equation}\label{eq:isotropic-gauge}
 \mathcal O'=\mathcal O+h,\qquad
 A_i'=A_i+h+\pi_i,\qquad \pi_i\in\operatorname{Per}(\mathcal O).
\end{equation}
The relative registrations are recovered modulo the same full period
group as in \eqref{eq:registration-coset}.
\end{theorem}
\begin{proof}
The preceding lemma and homogeneity of perimeter give
\eqref{eq:isotropic-ratio}; the first deficit is positive because
the first footprint has nonempty interior. The global occupation
inverse supplies $\mathcal P_i$ without knowledge of the footprint.
Steiner additivity gives
\[
 L_i(u)=\inf_C p_C^\circ(u)+\rho_i p_{-A_0}(u).
\]
Subtracting at two distinct scales proves
\eqref{eq:isotropic-footprint}, including its physical support
interpretation. No separate convexity of the envelopes or attainment
of their infima is used. Cancelling the centered footprint component
by component gives $\mathcal O-b_i$.

Equality of observed mean pairs for two experiments therefore implies
$\mathcal O-b_i=\mathcal O'-b_i'$ at every setting. Taking
$h=b_1'-b_1$ gives $b_i'-b_i-h\in\operatorname{Per}(\mathcal O)$
and proves necessity of \eqref{eq:isotropic-gauge}. Conversely,
translate each launch density by $h+\pi_i$. Translation of the table
by $h$ and then of a flight by a period preserve each collision bit,
for every direction. Thus the gauge is realized by identical binary
laws, and is sufficient as well.
\end{proof}

The perimeter normalization uses the raw forward mean, not merely its
difference with the reverse mean. The general occupation and period
inverse uses only that difference. Exact homothety and labelled
settings remain assumptions, although their origins and ratios are
recovered. No assertion about finite isotropic boundary minimax rates
is substituted for the proved compass and short-command rates.
